\documentclass[12pt]{article}

% --------------------
% PAQUETES BÁSICOS
% --------------------
\usepackage[spanish,es-noshorthands]{babel}
\usepackage[utf8]{inputenc}      % Si usas pdflatex
\usepackage[T1]{fontenc}
\usepackage{lmodern}

\usepackage{amsmath, amssymb, amsthm}
\usepackage{geometry}
\usepackage{booktabs}   
\usepackage{setspace}
\usepackage{graphicx}
\usepackage{tikz}
\usepackage{tikz-cd}
\usepackage[
  colorlinks=true,
  linkcolor=blue,
  citecolor=blue,
  urlcolor=red
]{hyperref}

\newtheorem{thm}{Teorema}[section]
\newtheorem{dfn}[thm]{Definición}
\newtheorem{lem}[thm]{Lema}
\newtheorem{ej}[thm]{Ejemplo}
\newtheorem{cor}[thm]{Corolario}
\newtheorem{obs}[thm]{Observación}
\newtheorem{con}[thm]{Conjetura}
\newtheorem{prop}[thm]{Proposición}

\DeclareMathOperator{\Frm}{Frm}
\DeclareMathOperator{\Ord}{Ord}
\DeclareMathOperator{\Top}{Top}
\DeclareMathOperator{\Sp}{sp}
\DeclareMathOperator{\pt}{pt}
\DeclareMathOperator{\id}{id}
\DeclareMathOperator{\tp}{tp}
\DeclareMathOperator{\fd}{fd}
% --------------------
% CONFIGURACIÓN
% --------------------
\geometry{margin=2.5cm}
\onehalfspacing

% --------------------
% DATOS DEL DOCUMENTO
% --------------------
\title{Reporte mensual de actividades}
\date{1-abril-2026 a 30-abril-2026}
\author{Juan Carlos Monter Cortés}

% --------------------
% INICIO DEL DOCUMENTO
% --------------------
\begin{document}

\maketitle
% \tableofcontents

\section{Actividades realizadas}
En \cite{simmons2004vietoris} y \cite{sexton2006ordinal} utilizan los espacios apilados y fuertemente apilados para estudiar la complejidad de los intervalos de admisibilidad. 
De manera particular, analizan la relación de esta clase de espacios con la modificación de Vietoris y los $V$-puntos. Los resultados que se presentan en este reporte
corresponden a las versiones libres de puntos de algunas de las caracterizaciones de los artículos anteriores.\\

Si consideramos $A\in \Frm$ y $(F, Q, \nabla)$ una terna HM tenemos el diagrama
\begin{equation}\label{Qlaxo}
\begin{tikzcd}
	A & A \\
	{\mathcal{O}S} & {\mathcal{O}S}
	\arrow["{v_F}", from=1-1, to=1-2]
	\arrow["{\mathcal{U}}"', from=1-1, to=2-1]
	\arrow["{\mathcal{U}}", from=1-2, to=2-2]
	\arrow["{v_\nabla}"', from=2-1, to=2-2]
\end{tikzcd}
\end{equation}
y se cumple que $\mathcal{U}\circ v_F\leq v_\nabla \circ \mathcal{U}$, donde $v_F, v_\nabla$ son los núcleos ajustados a los filtros abiertos 
$F, \nabla$, respectivamente, y $A_F$, $A_\nabla$ son los cocientes generados por estos núcleos. Queremos estudiar bajo qué condiciones el 
diagrama \eqref{Qlaxo} es conmutativo, es decir, $\mathcal{U}\circ v_F=v_\nabla \circ \mathcal{U}$.\\

En general, para cualquier terna $(F, Q, \nabla)$ se cumple que
\[
v_F\leq \mathcal{U}_*\circ v_\nabla\circ \mathcal{U}\leq \mathcal{U}_*\circ [Q']\circ\mathcal{U}
\]

Además, si $A$ es fuertemente apilado, entonces $v_F=\mathcal{U}_*\circ [Q']\circ \mathcal{U}$ y por lo tanto, si $A$ es fuertemente apilado, se cumple que 
\[
v_F=\mathcal{U}_*\circ v_\nabla\circ \mathcal{U}\Rightarrow \mathcal{U}\circ v_F=v_\nabla\circ \mathcal{U}
\]
donde la implicación es consecuencia de la fórmula del adjunto derecho.\\

Con esto hemos probado que si $A$ es fuertemente apilado, entonces el diagrama \eqref{Qlaxo} es conmutativo. El siguiente resultado establece
bajo qué condiciones tenemos que el recíproco es cierto.

\begin{prop}\label{EficienteFapilado}
  Si $A\in \Frm$ es eficiente y el diagrama 

\[\begin{tikzcd}
	A & A \\
	{\mathcal{O}S} & {\mathcal{O}S}
	\arrow["{v_F}", from=1-1, to=1-2]
	\arrow["{\mathcal{U}}"', from=1-1, to=2-1]
	\arrow["{\mathcal{U}}", from=1-2, to=2-2]
	\arrow["{v_\nabla}"', from=2-1, to=2-2]
\end{tikzcd}\]
conmuta, entonces $A$ es fuertemente apilado.
\end{prop}

\begin{proof}
Sean $A\in\Frm$, $S=\pt A$ y $(F, Q)$ un par HM. Por hipótesis, $A$ es eficiente y por resultados previos tenemos las implicaciones
\[
A\text{ eficiente }\Rightarrow \mathcal{O}S\text{ eficiente }\Rightarrow S \text{ apilado y empaquetado.}
\]
De aquí que, por empaquetado, si $Q\in \mathcal{Q}S$, entonces $Q\in \mathcal{C}S$ y $Q'\in \mathcal{O}S$. Luego, al ser $S$ apilado tenemos que
$\overline{Q}=\partial_\nabla^\infty(S)$, es decir, $Q'=v_\nabla(\emptyset)$ y al ser $\mathcal{O}S$ eficiente, $v_\nabla=u_{Q'}=[Q']$.\\
  
Además, si $v_\nabla\circ \mathcal{U}=\mathcal{U}\circ v_F$, entonces
\[
v_\nabla\circ \mathcal{U}\leq \mathcal{U}\circ v_F\quad \text{ y } \mathcal{U}\circ v_F\leq v_\nabla\circ \mathcal{U},
\]
es decir,
\[
[Q']\circ \mathcal{U}\leq \mathcal{U}\circ v_F\quad \text{ y } \mathcal{U}\circ v_F\leq [Q']\circ \mathcal{U}.
\]
Por propiedades del adjunto derecho 
\[
\mathcal{U}_*\circ [Q']\circ \mathcal{U}\leq v_F\quad \text{ y } v_F\leq\mathcal{U}_*\circ [Q']\circ \mathcal{U}.
\]
Por lo tanto $v_F=\mathcal{U}_*\circ [Q']\circ \mathcal{U}$, es decir, $A$ es fuertemente apilado.
\end{proof}

Un corolario que surge de la prueba anterior es el siguiente

\begin{cor}
Si $\mathcal{O}S$ es eficiente, entonces $S$ es fuertemente apilado.
\end{cor}

\begin{proof}
Ya que $S$ es fuertemente apilado si y solo si $v_\nabla=[Q']$, el resultado se sigue del hecho de que 
\[
\mathcal{O}S \text{ eficiente } \Rightarrow S \text{ apilado y empaquetado}
\] 
\end{proof}

Si $F\in A^\wedge$, podemos construir el conjunto
\[
M=\{m\in A\setminus F\mid a\leq m, \forall a\in A\setminus F\}.
\]
De manera particular, si $p\in \pt A$, entonces para el par HM $(F_p, \uparrow p)$ el conjunto $M$ es de la forma
\[
M=\{m\in A\setminus F_p\mid a\leq m, \forall a\in A\setminus F_p\}.
\]

Notemos que $A\setminus F_p=\{x\in A\mid x\leq p\}=\downarrow p$. Luego,
\[
M=\{m\in \downarrow p\mid a\leq m, \forall a\in \downarrow p\}=\{p\}.
\]

De aquí que $w_{F_p}(x)=\bigwedge\{p\mid x\leq p\}$, es decir, 
\begin{equation}\label{igualdadwF_p}
w_{F_p}(x)= \left\{ \begin{array}{lcc} 1, & \text{si}  &  x\nleq p\\ \\ 
	p, & \text{si} & x\leq p 
\end{array} \right.=w_p(x)
\end{equation}
para todo $x\in A$. Por lo tanto $w_{F_p}=w_p$.\\

\begin{prop}\label{maximosenpt}
Sea $A\in \Frm$ y $(F, Q)$ cualquier par HM. Si $v_F(0)=w_F(0)$, entonces todo punto de $A$ es máximo en $\pt A$.
\end{prop}

\begin{proof}
Consideremos $p,q\in \pt A$ tales que $p\leq q$. Notemos que para $q\in \pt A$, tenemos el par HM $(F_q, \uparrow q)$, donde 
\[
F_q=\{x\in A\mid \uparrow q\subseteq \mathcal{U}(x)\}=\{x\in A\mid (\forall r\in \pt A)[r\leq q\Rightarrow x\nleq q]\}.
\]

Por hipótesis, $v_{F_q}(0)=w_{F_q}(0)$ y por \eqref{igualdadwF_p}
\[
q=w_q(0)=w_{F_q}(0)=v_{F_q}(0)\leq v_{F_q}(p)=p,
\]
donde la última igualdad es consecuencia de que $p\notin F_q$. Por lo tanto, $p=q$ y todo punto es máximo en $\pt A$.
\end{proof}

Recordemos que para el morfismo reflexión espacial $\mathcal{U}_A$ tenemos el núcleo $\Sp_A$ dado por
\[
\Sp_A(x)=\bigwedge\{p\in \pt A\mid x\leq p\}
\]
para $x\in A$. Combinando este hecho con lo mencionado en la Proposición \ref{maximosenpt} podemos dar información sobre el colapso de 
los intervalos de admisibilidad.

\begin{prop}\label{focales}
Para $A\in \Frm$ y $(F, Q)$ cualquier par HM, se cumplen las siguientes afirmaciones:
\begin{enumerate}
\item Si $A$ es $T_1$ y apilado, entonces $v_F(0)=w_F(0)$.
\item Si $v_F(0)=w_F(0)$, entonces $A$ es apilado.
\item Si $v_F(0)=w_F(0)$ y el conjunto $Y=\{q\in \pt A\mid a\leq q\}$ es distinto de vacío para todo $a\in A$, entonces $A$ es $T_1$.
\end{enumerate}
\end{prop}

\begin{proof}
	\mbox{}

	\begin{enumerate}
		\item Si $A$ es $T_1$, se cumple que
		\[
		\mathcal{O}S\text{ es }T_1\Rightarrow S\text{ es }T_1\Rightarrow Q=M
		\]
		y al ser $A$ apilado
		\[
		v_F(0)=(\mathcal{U}_*\circ [Q']\circ \mathcal{U})(0)=(\mathcal{U}_*\circ[M']\circ \mathcal{U})(0)=w_F(0).
		\]
		\item Recordemos que para todo par HM se cumple 
		\begin{equation}\label{interad}
		v_F\leq \mathcal{U}_*\circ [Q']\circ \mathcal{U}\leq w_F
		\end{equation}
		y por hipótesis $v_F(0)=w_F(0)$, es decir, $v_F(0)=(\mathcal{U}_*\circ [Q']\circ \mathcal{U})(0)=w_F(0)$. Por lo tanto, $A$ es apilado.
		\item Consideremos $p\in \pt A$ y $a\in A$ tales que $p\leq a$. Por hipótesis $\{q\in \pt A\mid a\leq q\}\neq \emptyset$. Luego 
		\[
		p\leq a\leq \Sp (a)=\bigwedge{q\in \pt A\mid a\leq q}\leq q
		\]
		para algún $q\in Y$. Así, $p\leq a\leq q$ y por la Proposición \ref{maximosenpt}, $p=q$. Por lo tanto, $p=a$ y todo punto es máximo en $A$, 
		es decir, $A$ es $T_1$.
	\end{enumerate}
\end{proof}

Como consecuencia del resultado anterior tenemos el siguiente.

\begin{cor}\label{vF=wF}
Para $A\in \Frm$ y $(F, Q)$ cualquier par HM, se cumplen las siguientes afirmaciones::
\begin{enumerate}
\item Si $A$ es $T_1$ y fuertemente apilado, entonces $v_F=w_F$.
\item Si $v_F=w_F$, entonces $A$ es fuertemente apilado.
\item Si $v_F=w_F$ y el conjunto $Y=\{q\in \pt A\mid a\leq q\}$ es distinto de vacío para todo $a\in A$, entonces $A$ es $T_1$.
\end{enumerate}
\end{cor}

\begin{proof}
	\mbox{}

	\begin{enumerate}
	\item Si $A$ es fuertemente apilado $v_F=\mathcal{U}_*\circ [Q']\circ \mathcal{U}$ y por $T_1$ se cumple que $Q=M$. Por lo tanto $v_F=w_F$.
	\item Se sigue de \eqref{interad}.
	\item Se sigue de la Proposición \ref{focales}.
	\end{enumerate}
\end{proof}

Notemos que si para cada par HM $(F, Q)$ se cumple que $v_F(0)=w_F(0)$, entonces $A$ es apilado y todo punto es máximo en $\pt A$.
En particular, para cada $p\in \pt A$ y el par $(F_p, \uparrow p)$ se cumple que $v_{F_p}(0)=p$. Además $v_{F_p}(p)=v_{F_p}(0)=0_{A_{F_p}}$. \\

Si $A$ es $T_1$, $p$ es máximo en $A$, entonces $p$ es máximo en $A_{F_p}$ y por lo tanto, $A_{F_p}=\{q, 1\}$.\\

Ahora, $F_{F_p}=\nabla(v_{F_p})$ y $x\in F_p$ si y solo si $v_{F_p}(x)=1$. De esta manera, bajo las hipótesis anteriores, se cumple que 
\[
v_{F_p}(x)= \left\{ \begin{array}{lcc} 1, & \text{si}  &  x\nleq p\\ \\ p, & \text{si} & x\leq p 
\end{array} \right.=  \left\{ \begin{array}{lcc} 1, & \text{si}  &  x\in F_p\\ \\ p, & \text{si} & x\not\in F_p \end{array} \right.=w_{F_p}(x).
\]

Por lo tanto, si $A$ es $T_1$ y apilado, entonces $v_{F_p}=w_{F_p}$ para cada $p\in \pt A$. Con esto en mente, podríamos preguntarnos si existe alguna manera de 
recuperar la noción fuertemente apilado por medio de $T_1$ y apilado.\\

Responder esto último y verificar si los diferentes resultados de \cite{simmons2004vietoris} y \cite{sexton2006ordinal} en los que aparecen las nociones de
apilado y fuertemente apilado se comportan de manera similar cuando los trasladamos al contexto de marcos, son actividades que realizarán en las próximas semanas.

\bibliographystyle{amsalpha}
\bibliography{research2}


\end{document}